\chapter{The twelve-sector modular data}
\label{ch:modular}

Modular data translate categorical information into formulas for
characters, defects, and boundaries.  They are powerful because a
small pair of matrices constrains infinitely many states.  They are
also easy to overinterpret: modular matrices do not determine all
operator products, and matching them does not in general prove an
equivalence of braided categories.  This chapter derives the main
Haagerup formulas and works through a nontrivial Verlinde calculation.
It ends by constructing the categorical algebra underlying the
selected holomorphic extension.

\section{Tori, twists, and the central-charge phase}

A torus with complex parameter $\tau$ can be described using different
bases of its two cycles.  The basic changes are
\[
 \tau\longmapsto-1/\tau,\qquad \tau\longmapsto\tau+1.
\]
For a rational chiral algebra with irreducible modules $M_i$, its
characters
\[
 \chi_i(\tau)=\Tr_{M_i}q^{L_0-c/24},\qquad q=e^{2\pi i\tau},
\]
transform by matrices $S$ and $T_{\rm ch}$.  If the lowest weight
of $M_i$ is $h_i$, then
\begin{equation}
 (T_{\rm ch})_{ii}=e^{2\pi i(h_i-c/24)}.
 \label{eq:character-T-general}
\end{equation}
The corresponding categorical twist is
\begin{equation}
 \theta_i=e^{2\pi i h_i},\qquad
 \Theta=\operatorname{diag}(\theta_i),\qquad
 T_{\rm ch}=e^{-2\pi i c/24}\Theta.
 \label{eq:twist-character-distinction}
\end{equation}
The phase multiplying $\Theta$ is not optional.  Forgetting it
confuses a topological Dehn twist with the action on a chiral
character including its vacuum energy.

The normalized $S$-matrix can be obtained from a Hopf-link trace.
We choose its orientation to agree with the finite Fourier convention
\eqref{eq:Z3-center-S}.  In a unitary category this is the normalized
trace of inverse double braiding.  Reversing all orientation
conventions gives equivalent formulas after the appropriate changes
to twists and modular generators.  The Haagerup $S$-matrix below is
real, so complex conjugation of its entries has no effect.

Let
\[
 \mathcal D=\left(\sum_i d_i^2\right)^{1/2}
\]
be the total quantum dimension of a unitary modular category.
Normalization gives
\begin{equation}
 S_{00}=\mathcal D^{-1},\qquad S_{0i}=d_i/\mathcal D.
 \label{eq:S-vacuum-row}
\end{equation}
The vacuum row is therefore positive.  Charge conjugation is the
matrix $C_{ij}=\delta_{i,j^*}$, and
\begin{equation}
 S^2=C,\qquad
 (S\Theta)^3=\xi C,
 \qquad
 \xi=\frac1{\mathcal D}\sum_i d_i^2\theta_i.
 \label{eq:modular-relations-phase}
\end{equation}
For the representation category of a suitable unitary rational
chiral theory, $\xi=e^{2\pi i c/8}$.  Substituting
\eqref{eq:twist-character-distinction} removes that phase and gives
$(ST_{\rm ch})^3=C$.

The twist determines $h_i$ only modulo integers.  Knowing
$\theta=e^{4\pi i/13}$ permits lowest weights $2/13$, $15/13$,
$28/13$, and so on.  Positivity, branching, and actual character
coefficients are needed to decide which lift occurs.

\section{The Haagerup labels and matrix blocks}
\index{modular data!Haagerup center}

Use the ordering
\begin{equation}
 (0,b,a,c_0,c_1,c_2,d_1,d_2,d_3,d_4,d_5,d_6).
 \label{eq:H3-center-order}
\end{equation}
The first label is the vacuum. The $d_j$ are names of sectors;
their quantum dimensions will all be $3d$. Write $r=\sqrt{13}$ and set
\[
 A=\frac16-\frac1{2r},\qquad
 B=\frac16+\frac1{2r},\qquad
 u=\frac13,\qquad v=\frac1r.
\]
Let $\mathbf1_n$ denote the column of $n$ ones. In blocks of
sizes $1,1,4,6$, the matrix is
\begin{equation}
 S=\begin{pmatrix}
 A&B&u\mathbf1_4^{\mathsf t}&v\mathbf1_6^{\mathsf t}\\
 B&A&u\mathbf1_4^{\mathsf t}&-v\mathbf1_6^{\mathsf t}\\
 u\mathbf1_4&u\mathbf1_4&C_4&0\\
 v\mathbf1_6&-v\mathbf1_6&0&D_6
 \end{pmatrix},
 \label{eq:H3-modular-S}
\end{equation}
where
\begin{equation}
 C_4=\frac13\begin{pmatrix}
 2&-1&-1&-1\\
 -1&2&-1&-1\\
 -1&-1&-1&2\\
 -1&-1&2&-1
 \end{pmatrix},\qquad
 (D_6)_{ij}=-\frac2r\cos\frac{2\pi ij}{13}.
 \label{eq:H3-modular-blocks}
\end{equation}
The twist matrix is
\begin{align}
 \Theta=\operatorname{diag}\bigl(&1,1,1,1,
 e^{2\pi i/3},e^{4\pi i/3},
 e^{2\pi i6/13},e^{2\pi i11/13},e^{2\pi i2/13},\notag\\
 &e^{2\pi i5/13},e^{2\pi i7/13},e^{2\pi i8/13}\bigr).
 \label{eq:H3-modular-twists}
\end{align}
Equivalently, $\theta_{d_j}=e^{2\pi i6j^2/13}$ for
$1\leq j\leq6$.  The six exponents are the quadratic nonresidues
modulo thirteen.

\begin{constructioninput}[The Haagerup center data]
The standard unitary Haagerup center has the modular matrices
\eqref{eq:H3-modular-S}--\eqref{eq:H3-modular-twists} and the
forgetful data of \eqref{eq:H3-forgetful}.  In the CFT construction,
identifying the actual module category with this center requires
the separate structural and analytic arguments developed later.
Writing down these matrices is not a substitute for those arguments.
\end{constructioninput}

In the bootstrap notation frequently used for these sectors, the
same order is
\[
 (1,\pi_1,\pi_2,\sigma_1,\sigma_2,\sigma_3,
 \mu_5,\mu_2,\mu_1,\mu_3,\mu_6,\mu_4).
\]
The permutation among the final six labels is fixed by their twists.
Keeping equal numerical subscripts on $d_j$ and $\mu_j$ would
misidentify sectors.  We retain \eqref{eq:H3-center-order} in
calculations and state any translation when comparing conventions.

\section{Checking normalization without a computer}

The first entry can be written
\[
 A=\frac{\sqrt{13}-3}{6\sqrt{13}}
 =\frac1{3s}=\frac1{\mathcal D}.
\]
Consequently the ratios of the vacuum-row entries to $A$ are
\begin{equation}
 (\qdim(M_i))_i=(1,d^2,s,s,s,s,3d,3d,3d,3d,3d,3d).
 \label{eq:H3-central-dimensions}
\end{equation}
For example $u/A=s$, while $v/A=3d$.  The total squared dimension
is indeed
\begin{align*}
 1+d^4+4s^2+6(3d)^2
 &=1+d^4+4(1+d^2)^2+54d^2\\
 &=5+5d^4+62d^2\\
 &=9(1+d^2)^2=(3s)^2.
\end{align*}
The last step follows from $d^4=11d^2-1$, itself a consequence
of $d^2=3d+1$.

We can also verify $S^2=I$ using finite Fourier sums.  For
$m\not\equiv0\pmod{13}$,
\begin{equation}
 \sum_{j=1}^6\cos\frac{2\pi mj}{13}=-\frac12.
 \label{eq:cos13-sum}
\end{equation}
It is half the sum over all nonzero residues, which equals $-1$.
Using $2\cos x\cos y=\cos(x+y)+\cos(x-y)$ gives
\begin{equation}
 \sum_{j=1}^6
 \cos\frac{2\pi ij}{13}\cos\frac{2\pi kj}{13}
 =\frac{13}{4}\delta_{ik}-\frac12
 \quad(1\leq i,k\leq6).
 \label{eq:cos13-orthogonality}
\end{equation}
Writing $I_n$ for the identity matrix and $J_n$ for the $n$-by-$n$
matrix of ones, we obtain
\[
 D_6^2=I_6-\frac2{13}J_6,
 \qquad D_6\mathbf1_6=v\mathbf1_6.
\]
The first two columns of $S$ contribute $2v^2J_6$, exactly the
missing term in the bottom-right block of $S^2$.

For the middle block, let $P$ fix its first two coordinates and
exchange its last two.  Then $C_4=P-J_4/3$, so
\[
 C_4^2=I_4-\frac29J_4,
 \qquad C_4\mathbf1_4=-\frac13\mathbf1_4.
\]
The first two columns contribute $2u^2J_4$, again exactly the
missing term.  Mixed blocks vanish by these row sums and
$A+B=u$, $A-B=-v$.  Finally,
\[
 A^2+B^2+4u^2+6v^2=1,
 \qquad 2AB+4u^2-6v^2=0.
\]
This verifies every block of $S^2=I$.  All twelve simples are
therefore self-dual.  This statement concerns the center labels;
the six-object category still has $\alpha^*=\alpha^2$.
For a simple object, ribbon duality gives $\theta_{x^*}=\theta_x$:
the dual morphism is a linear transpose. Taking a Hermitian adjoint
instead conjugates the scalar. Thus the self-dual label $c_1$ can
have the nonreal twist $e^{2\pi i/3}$; categorical duality and
right-moving complex conjugation are different operations.

\section{A Gauss sum and central charge modulo eight}
\index{central charge!modulo eight}

Let $\zeta=e^{2\pi i/13}$.  The standard quadratic Gauss-sum
evaluation gives
\begin{equation}
 G_-:=\sum_{j=1}^6\zeta^{6j^2}
 =\frac{-1-\sqrt{13}}2.
 \label{eq:nonresidue-Gauss13}
\end{equation}
One way to see its form is to separate the twelve nonzero residues
into squares and nonsquares.  Their two sums add to $-1$, while
their difference is the quadratic Gauss sum $\sqrt{13}$ with
the standard positive choice for this prime.  The coefficient $6$
selects the nonsquares.

The Gauss--Milgram numerator is now
\begin{align}
 p^+&=1+d^4+s^2(1+1+\omega+\omega^2)
       +9d^2G_-\notag\\
 &=1+d^4+s^2+\frac{9d^2}{2}(-1-\sqrt{13})
 =3s=\mathcal D.
 \label{eq:H3-Gauss-Milgram}
\end{align}
Thus $\xi=1$.  This also follows abstractly because a Drinfeld
center has trivial multiplicative central charge.

If a rational chiral algebra has representation category
$Z(\HH_3)$, this constrains its chiral central charge to
$c\equiv0\pmod8$.  The construction in this book realizes the
choice $c=8$.

There is no contradiction with a proposed full CFT of central charge
near two carrying Haagerup \emph{defect symmetry}.  The latter claim
does not identify its chiral representation category with
$Z(\HH_3)$.  A symmetry category acts on a full theory, while
the modular category of chiral modules describes a chosen rational
chiral algebra.  These are different pieces of structure.  Likewise,
the holomorphic $E_8$ chiral algebra has a trivial representation
category while its full theory can possess nontrivial defects
preserving a smaller algebra $U$.

\section{Verlinde diagonalizes fusion}
\index{Verlinde formula}

The fusion matrices of a modular category commute and are
simultaneously diagonalized by $S$.  In our convention their
structure constants are
\begin{equation}
 N_{ij}^{\ k}=\sum_\ell
 \frac{S_{i\ell}S_{j\ell}\overline{S_{k\ell}}}
 {S_{0\ell}}.
 \label{eq:Verlinde-general}
\end{equation}
The eigenvalue of the line $i$ in sector $\ell$ is
$S_{i\ell}/S_{0\ell}$.  The vacuum eigenvalue is $d_i$,
recovering the dimension character.

For Haagerup all entries are real, so the bar in
\eqref{eq:Verlinde-general} can be omitted.  Although individual
entries contain radicals and cosines, the sum gives nonnegative
integers.  That integrality is a strong consistency condition, but
its validity alone would still not reconstruct an associator or
prove that a candidate modular table comes from a CFT.

\begin{example}[The identity coefficient]
For any two labels,
\[
 N_{ij}^{\ 0}=\sum_\ell S_{i\ell}S_{j\ell}
 =(S^2)_{ij}=\delta_{ij}.
\]
This agrees with self-duality of every Haagerup center simple.
The vacuum fusion rule follows similarly from unitarity:
$N_{0i}^{\ k}=\delta_{ik}$.
\end{example}

\section{A full nontrivial fusion product from cosine sums}

We compute $d_i\otimes d_i$.  Here $d_i$ is a center label,
whereas the italic scalar $d$ remains the quantum dimension of a
noninvertible Haagerup line.  Write
$c_m(j)=\cos(2\pi mj/13)$.

For an output $d_k$, the first two columns of Verlinde contribute
\[
 v^3\left(\frac1A-\frac1B\right)=\frac9{13},
\]
using $AB=1/117$.  The six final columns contribute
$-(8/13)\sum_j c_i(j)^2c_k(j)$.  Thus
\begin{equation}
 N_{d_i d_i}^{\ \ d_k}
 =\frac9{13}-\frac8{13}
 \sum_{j=1}^6 c_i(j)^2c_k(j).
 \label{eq:Verlinde-di-square-sum}
\end{equation}
Product-to-sum gives
\[
 c_i^2c_k=\frac12c_k+\frac14c_{2i+k}+\frac14c_{2i-k}.
\]
Let $[m]$ denote the representative of $\{m,-m\}$ modulo
thirteen in $\{1,\ldots,6\}$.  If $k\ne[2i]$, all three
frequencies are nonzero modulo thirteen, so
\eqref{eq:cos13-sum} makes the sum $-1/2$.  Equation
\eqref{eq:Verlinde-di-square-sum} then gives one.
If $k=[2i]$, precisely one of $2i+k$ and $2i-k$ vanishes;
its cosine sum is six instead of $-1/2$.  The triple-product sum
becomes $9/8$, giving zero.  Therefore
\begin{equation}
 N_{d_i d_i}^{\ \ d_k}=1-\delta_{k,[2i]}.
 \label{eq:Verlinde-di-square-residual}
\end{equation}

For any of $a,c_0,c_1,c_2$, only the first two columns contribute,
and their sum is
\[
 uv^2\left(\frac1A+\frac1B\right)=1.
\]
The coefficient of $b$ is zero: its first two columns contribute
$11/13$, while the final block contributes
$-\sum_j(D_6)_{ij}^2=-11/13$.  The vacuum coefficient is one.
Putting everything together,
\begin{equation}
 d_i^2=\one+a+c_0+c_1+c_2+
 \sum_{\substack{1\leq k\leq6\\k\ne[2i]}}d_k.
 \label{eq:Verlinde-di-square}
\end{equation}
Here $\one=M_0$ denotes the vacuum object, so the first term has
dimension one.  The dimension check is
\[
 (3d)^2=1+4s+5(3d)=9+27d,
\]
exactly as required by $d^2=1+3d$.

Another compact product is
\[
 b^2=\one+b+a+c_0+c_1+c_2+\sum_{j=1}^6d_j.
\]
Its dimension identity is $d^4=1+d^2+4s+18d$.
These examples show explicitly how finite Fourier cancellation turns
algebraic and trigonometric entries into integer fusion rules.

\section{The canonical extension algebra}
\index{Lagrangian algebra}

Let $F:Z(\HH_3)\to\HH_3$ be the forgetful functor and $I$
its right adjoint.  The object $I(\one)$ carries a canonical
connected commutative separable algebra structure.  Its summands
follow from adjunction:
\[
 \Hom_{Z(\HH_3)}(\mu,I(\one))
 \simeq\Hom_{\HH_3}(F(\mu),\one).
\]
Using the first row of the forgetful table gives
\begin{equation}
 \mathcal A:=I(\one)\simeq\one\oplus b\oplus a^{\oplus2}.
 \label{eq:canonical-H3-algebra}
\end{equation}
Here $\one=M_0$ is the vacuum center object; the numerical label
$0$ remains its index in the modular matrices.  Its quantum dimension is
\begin{equation}
 \qdim(\mathcal A)=1+d^2+2s=3s=\mathcal D.
 \label{eq:canonical-H3-algebra-dimension}
\end{equation}
Thus $(\qdim\mathcal A)^2$ equals the global dimension of the
center.  Such an algebra is called Lagrangian.

For a connected commutative separable algebra in a nondegenerate
braided category, the dimension formulas are
\begin{align}
 \operatorname{Dim}(\mathcal M_{\mathcal A})
 &=\frac{\operatorname{Dim}(\mathcal M)}{\qdim\mathcal A},\notag\\
 \operatorname{Dim}(\mathcal M_{\mathcal A}^{\rm loc})
 &=\frac{\operatorname{Dim}(\mathcal M)}{(\qdim\mathcal A)^2}.
 \label{eq:algebra-module-dimensions}
\end{align}
For $\mathcal M=Z(\HH_3)$ these become $\mathcal D$ and one,
respectively.  More strongly, the canonical algebra has
\[
 Z(\HH_3)_{\mathcal A}\simeq\HH_3,
 \qquad Z(\HH_3)_{\mathcal A}^{\rm loc}\simeq\operatorname{Vec}.
\]
The first equivalence retains all modules; the second keeps only
local ones.

\subsection{Why a Lagrangian algebra recovers a center}
\label{sec:lagrangian-reconstruction}
\index{Lagrangian algebra!center reconstruction}
\index{quantum dimension!module normalization}
Let $A$ be a connected commutative separable algebra in the
nondegenerate braided category $\mathcal M$.
The tensor unit of the module category $\mathcal M_A$ is $A$ itself.
Its dimension there must be one even when its dimension in
$\mathcal M$ is large. The normalization is
\[
 \FPdim_{\mathcal M_A}(M)
   =\frac{\FPdim_{\mathcal M}(M)}{\FPdim_{\mathcal M}(A)},
\]
where the numerator forgets the module action. In our positive
spherical setting these are also the categorical dimensions.
For instance, in $Z(\operatorname{Vec}_{\Z_3})$ take
$A=\bigoplus_k(0,k)$. The flux module
$x_g=\bigoplus_k(g,k)$ has ambient dimension three and module
dimension one. The three flux modules consequently have global
dimension three, as the dimension formula predicts. Their actions
and relative tensor products are worked out in
Chapter~\ref{ch:defects}.

\index{central functor}
There is also an actual functor behind reconstruction. The free
module functor is
$G_A:\mathcal M\to\mathcal M_A$, $G_A(X)=X\otimes A$.
It is \emph{central}: the ambient braiding supplies coherent
half-braidings of $G_A(X)$ with every $A$-module. Equivalently,
it lifts to a braided tensor functor
\[
 \widetilde G_A:\mathcal M\longrightarrow Z(\mathcal M_A).
\]
This is the meaning of a central functor; it need not take values in
a braided category before the lift.
For a fusion category, $\FPdim(\mathcal M)$ means the sum of the
squares of its simple Frobenius--Perron dimensions. In the unitary
setting it agrees with the global dimension $\operatorname{Dim}(\mathcal M)$.

\begin{theorem}[Lagrangian reconstruction]
\label{thm:lagrangian-reconstruction}
Let $\mathcal M$ be a nondegenerate braided fusion category and
$A$ a connected commutative separable algebra. If
$\FPdim(A)^2=\FPdim(\mathcal M)$, then the central lift gives
a braided equivalence
$\mathcal M\simeq Z(\mathcal M_A)$, carrying $A$ to the canonical
algebra $I(\one)$. Here $I$ is the right adjoint of the forgetful
functor $Z(\mathcal M_A)\to\mathcal M_A$.
\end{theorem}
The structural input is that this lift is fully faithful: it induces
an isomorphism on every morphism space. Once that is known, the
dimension formulas explain why the Lagrangian condition
makes it exhaustive:
\[
 \FPdim Z(\mathcal M_A)
 =\left(\frac{\FPdim\mathcal M}{\FPdim A}\right)^2
 =\FPdim\mathcal M.
\]
A fully faithful fusion subcategory with the full positive global
dimension cannot omit any simple object. This general theorem is
what turns a holomorphic extension $H/U$ into the identification
$\Rep(U)\simeq Z(\Rep(U)_H)$ later in the construction.

To apply the canonical Haagerup algebra $\mathcal A$ to $U$, one
needs an actual unitary ribbon
equivalence $\Rep(U)\simeq Z(\HH_3)$ and the relevant extension
theorem.  Transporting $\mathcal A$ then gives a selected holomorphic
extension.  Matching $S$ and $\Theta$ alone does not provide that
transport.  The original embedded extension in $V_{E_8}$ may identify
its boundary category only up to $\HH_2$ or $\HH_3$; selecting
the canonical $\HH_3$ algebra is a separate categorical choice.

\section{A modular-invariant vector and the number 248}

Let $m=(1,1,2,0,\ldots,0)^{\mathsf t}$ be the multiplicity vector
of \eqref{eq:canonical-H3-algebra}.  Direct substitution in
\eqref{eq:H3-modular-S} gives
\begin{equation}
 Sm=m,\qquad \Theta m=m.
 \label{eq:extension-modular-vector}
\end{equation}
For example the first component is
$A+B+2u=1$, the $a$ component is $2u+2(2/3)=2$,
and every $d_j$ component is $v-v=0$.
Therefore
\begin{equation}
 Z_{\rm ext}=|\chi_0+\chi_b+2\chi_a|^2
 \label{eq:extension-partition-function}
\end{equation}
is modular invariant whenever the characters transform by these
matrices.  This is a consistency consequence of the extension
algebra.  Conversely, the same invariant expression written for
formal character candidates would not construct local vertex
operators or prove sewing.

For the actual $c=8$ construction, the top spaces of $M_b$ and
$M_a$ have weight one and dimensions $80$ and $81$, while
$\dim U_1=6$.  Restricting the selected extension to $U$ gives
\[
 H'|_U\simeq U\oplus M_b\oplus2M_a.
\]
Consequently its weight-one space has dimension
\begin{equation}
 6+80+2\cdot81=248.
 \label{eq:248-branching-check}
\end{equation}
The equality is a valuable branching check.  Identifying the
extension with $V_{E_8}$ uses the holomorphic $c=8$ classification
and the extension's established hypotheses, not merely the familiar
number $248$.

For this $c=8$ theory,
$T_{\rm ch}=e^{-2\pi i/3}\Theta$.  The module called $d_3$
has actual lowest weight $15/13$, although its twist exponent is
$2/13$.  This one example captures why a modular table cannot be
read as a complete table of lowest energies.

\section{Exercises}

\begin{exercise}\label{ex:Z3-S-square}
Using \eqref{eq:Z3-center-S}, prove $S^2=C$ by performing the
two finite geometric sums.  Explain why replacing $C$ by the
identity would be incorrect for this example.

\par\smallskip\noindent{\footnotesize\normalfont
\hyperref[sol:ex:Z3-S-square]{Solution on page~\pageref*{sol:ex:Z3-S-square}.}}\par
\end{exercise}

\begin{exercise}\label{ex:H3-S-square}
Complete the block verification of $S^2=I$ for
\eqref{eq:H3-modular-S}, including the mixed $4\times6$ block.
Derive $AB=1/117$ and $A^2+B^2=11/117$.

\par\smallskip\noindent{\footnotesize\normalfont
\hyperref[sol:ex:H3-S-square]{Solution on page~\pageref*{sol:ex:H3-S-square}.}}\par
\end{exercise}

\begin{exercise}\label{ex:Gauss13}
List the six squares and six nonsquares modulo thirteen.  Verify
that $6j^2$, for $1\leq j\leq6$, runs through all nonsquares.
Use \eqref{eq:nonresidue-Gauss13} to check every algebraic step
of \eqref{eq:H3-Gauss-Milgram}.

\par\smallskip\noindent{\footnotesize\normalfont
\hyperref[sol:ex:Gauss13]{Solution on page~\pageref*{sol:ex:Gauss13}.}}\par
\end{exercise}

\begin{exercise}\label{ex:Verlinde-d1}
Write the product $d_1^2$ explicitly from
\eqref{eq:Verlinde-di-square}.  Which $d_j$ is absent?  Repeat
for $d_4^2$, and check dimensions in both cases.

\par\smallskip\noindent{\footnotesize\normalfont
\hyperref[sol:ex:Verlinde-d1]{Solution on page~\pageref*{sol:ex:Verlinde-d1}.}}\par
\end{exercise}

\begin{exercise}\label{ex:modular-extension}
Prove \eqref{eq:extension-modular-vector} component by component.
Show that the rank-one coefficient matrix $Z=mm^{\mathsf t}$
commutes with both $S$ and $\Theta$.  State two additional kinds
of information needed before interpreting this matrix as an actual
local CFT.

\par\smallskip\noindent{\footnotesize\normalfont
\hyperref[sol:ex:modular-extension]{Solution on page~\pageref*{sol:ex:modular-extension}.}}\par
\end{exercise}

\begin{exercise}\label{ex:central-charge-scope}
Explain why $\Rep(U)\simeq Z(\HH_3)$ at $c=8$ and a
conjectural full $c=2$ theory with $\HH_3$ symmetry are not
inconsistent statements.  Identify exactly which object carries
the central-charge-modulo-eight constraint in the first statement.

\par\smallskip\noindent{\footnotesize\normalfont
\hyperref[sol:ex:central-charge-scope]{Solution on page~\pageref*{sol:ex:central-charge-scope}.}}\par
\end{exercise}
